%========================================
%PREAMBOLO
%========================================

% === Impostazione del documento ==========================
\documentclass[12pt,a4paper,twoside,english,hidelinks]{book}
\pagenumbering{arabic}
\usepackage{setspace}
\onehalfspace
\usepackage[a-1b]{pdfx}
\usepackage{hyperref}
\usepackage{blindtext}

% === Regolazione dei margini =============================
\addtolength{\oddsidemargin}{30pt}
\addtolength{\evensidemargin}{-30pt}
\usepackage{fancyhdr}
\usepackage{multirow}
\usepackage{multicol}
\usepackage[section]{placeins}

% === Impostazione dei font ===============================
\usepackage[T1]{fontenc}
\usepackage[utf8]{inputenc}
%\usepackage[italian]{babel}
\usepackage[english]{babel}
\usepackage{ae}
\usepackage{relsize}
\usepackage{csquotes}
\usepackage{amsmath}
\usepackage{amsfonts}
\usepackage{mathdots}
\usepackage{mathtools}

\hypersetup{
%%pdfencoding=unicode
	bookmarksnumbered=true,
	linkcolor=black,
	citecolor=black,
	pagecolor=black,
	urlcolor=black,
%    pdfauthor={Argentieri Francesco}
}
\usepackage{verbatim}
\usepackage{alltt,footnote}
\DeclareMathOperator{\sgn}{sgn}
\DeclarePairedDelimiter{\abs}{\lvert}{\rvert}
\DeclarePairedDelimiter{\norma}{\lVert}{\rVert}

% === Integrazione delle figure ===========================
\usepackage{graphicx}
\graphicspath{{./imgs/}}
\renewcommand{\figurename}{Fig.}
\usepackage{subfig}
\usepackage{xcolor}

% === Per gli algoritmi ===================================
\usepackage{algorithmicx}
\usepackage[ruled]{algorithm}
\usepackage{algpseudocode}

% === Per le tabelle ======================================
\usepackage{tabularx}
%\usepackage[table]{xcolor}

%==Per le figure latex======================================
%Package e librerie per TikZ e PGF, le librerie non sono tutte necessarie a questo documento LATEX.
\usepackage{tikz,fp,ifthen,fullpage}
\usepackage{pgfmath}
\usetikzlibrary{decorations.pathmorphing,backgrounds,fit,calc,through}
\usetikzlibrary{arrows}
\usetikzlibrary{shapes,decorations,shadows}
\usetikzlibrary{fadings}
\usetikzlibrary{patterns}
\usetikzlibrary{mindmap}
\usetikzlibrary{decorations.text}
\usetikzlibrary{decorations.shapes}
\usepackage{pgfplots}
\usepgfplotslibrary{colorbrewer}
\usepgfplotslibrary{patchplots}
\pgfplotsset{compat=1.13}

%==Per APDL e linguaggi======================================
\usepackage{listings,color}
\definecolor{codeBlue}{RGB}{0,0,255}
\definecolor{codeGreen}{RGB}{0,128,0}
\definecolor{codePurple}{RGB}{128,0,255}
\definecolor{codeOrange}{RGB}{255,128,128}
\definecolor{codePink}{RGB}{255,0,255}
\definecolor{codeGray}{rgb}{0.5,0.5,0.5}
\lstdefinestyle{apdl-modified}
    {
        extendedchars=true,
        alsoletter={*},
        alsoletter={'},
        keywordstyle=\color{codeOrange},
        keywordstyle=[2]\color{codePurple},
        keywordstyle=[3],
        keywordstyle=[4]\color{codePink},
        otherkeywords={
            /solve,
            /angle
        },
        keywords=[2]{
            *abbr,
            *dim
        },
        keywords=[3]{
            m,
            d
        },
        keywords=[4]{
            ',
            -,
            ",
            \%,
            (,
            ),
            ,,
            .,
            :,
            ;,
            ?,
            ^,
            ~,
            +,
            <,
            =,
            >
        },
        sensitive=false, % keywords are not case-sensitive
        morecomment=[l]{!}, % l is for line comment
        commentstyle=\color{codeGray}, % style of comments
        numberstyle=\tiny\color{codeGray}
    }
\lstdefinelanguage{APDL}{style=apdl-modified}

% === Per la bibliografia multicolonna ====================
\usepackage{etoolbox}
\patchcmd{\thebibliography}{\list}{\begin{multicols}{2}\smaller\list}{}{}
\appto{\endthebibliography}{\end{multicols}}

%========================================
% TESTO DELLA TESI
%========================================
\begin{document}
	% === Frontespizio ====================================
	\pagestyle{empty}
	%%%%%%%%%%%%%%%%%%%%%%%%%%%%%%%%%%%%%%%%%%%%%%%%%%%%%%%%%%%
% Frontespizio
%%%%%%%%%%%%%%%%%%%%%%%%%%%%%%%%%%%%%%%%%%%%%%%%%%%%%%%%%%%
\begin{titlepage}
 \begin{center}
 \includegraphics[width=3.5cm]{unitn}\\
 \vspace{1.5em}
 {\Large \textsc{University of study of Trento}}\\
 \vspace{1.5em}
 {\Large \textsc{Department of Industrial Engineering}}\\
 \vspace{4em}
 {\normalsize Master of Science in Mechatronics Engineering}\\
 \vspace{1.5em}
 {\Large \textsc{Mechanical design and machine elements}}\\
 \vspace{4em}
 {\LARGE\textbf{
 	Report homework
 }}\\
 \end{center}

\vskip 2.0cm
\vskip 2.0cm
 \begin{center}
 \begin{tabular}{c c c c c c c c}
 Professor & & & & & & & Candidate \\[0.2cm]
 \large{Prof. Brunelli Davide} & & & & & & & \large{Francesco Argentieri}\\[0.4cm]
  & & & & & & & ID 183892\\[0.2cm]
% \large{Dott. Egidio Labbate}& & & & & & &\\
 \end{tabular}
 \end{center}

\vskip 1.5cm
\begin{center}
{\normalsize academic year 2018/2019}
\end{center}
\end{titlepage}

\clearpage{\pagestyle{empty}\cleardoublepage}

	% === Indice ==========================================
    \tableofcontents
    \listoffigures

	% === Capitoli Tesi ===================================
	\pagestyle{plain}
	%%%%%%%%%%%%%%%%%%%%%%%%%%%%%%%%%%%%%%%%%%%%%%%%%%%%%%%%%%%

\chapter{Introduzione}
\label{chap:Introduzione}

%%%%%%%%%%%%%%%%%%%%%%%%%%%%%%%%%%%%%%%%%%%%%%%%%%%%%%%%%%%

\section{Contesto applicativo}
\label{sec:intro1}

Ciao ragazzi :D questo è un template che potete utilizzare per scrivere la vostra tesi in \LaTeX (sì, il nome di questo linguaggio è tutto un programma...)

Cos'è \LaTeX ? Cercatevelo su Wikipedia\footnote{http://it.wikipedia.org/wiki/LaTeX}!

A parte gli scherzi... è un linguaggio che vi permette (in poche parole) di creare documenti accademici (e non) con uno stile molto professionale. Gran parte del lavoro sporco (creazione dei capitoli, delle sezioni, dell'indice, della bibliografia, gestione dei margini, ecc...) viene interamente gestito da \LaTeX , le poche cose da configurare sono già state impostate in questo template... (quindi in poche parole avete già tutto pronto, stronzi!)

In questo pdf è spiegato un po' come far funzionare il tutto, ovvero:
\begin{itemize}
\item Dove scaricare l'IDE e come configurarlo
\item Come scaricare il compilatore
\item Come iniziare a personalizzare questo template
\end{itemize}

Cercherò di utilizzare più elementi \LaTeX possibile nello scrivere questa guida (tabelle, elenchi puntati, footnote, immagini...) così quando andrete a leggere il codice vi imparate pure qualcosa, caproni (<3)

%%%%%%%%%%%%%%%%%%%%%%%%%%%%%%%%%%%%%%%%%%%%%%%%%%%%%%%%%%%

\section{Motivazioni e obiettivi}
\label{sec:intro2}

Il principale motivo che mi spinge a creare questo pdf è quello di risparmiarvi gran parte delle rotture che si trovano quando ci si avvicina al mondo \LaTeX... insomma mi auguro che questa guida vi permetta di avere un buon punto d'inizio.

Come già spiegato nella sezione precedente, \LaTeX offre tantissimi servizi utili ed uno stile professionale unico, cose che su altri programmi (come Microsoft Word) potete anche sognarvi.

Insomma... con \LaTeX potete presentare una tesi fatta come si deve :)

%%%%%%%%%%%%%%%%%%%%%%%%%%%%%%%%%%%%%%%%%%%%%%%%%%%%%%%%%%%

\section{Risultati raggiunti}
\label{sec:intro3}

Nella Figura \ref{fig:latexVsWord} potete ammirare quanto \LaTeX sia più figo di Microsoft Word, ooooooh...

\begin{figure}[h] %La "h" indica che la figura si posizionerà "here". Usate "t" per "top" e "b" per "bottom"
\centering %centrata
%\includegraphics[width=0.8\linewidth]{latexVsWord} %larghezza e nome del file
\caption{Oooooooh che figo \LaTeX, ooooooooh} %didascalia
\label{fig:latexVsWord} %label per il riferimento
\end{figure}

%%%%%%%%%%%%%%%%%%%%%%%%%%%%%%%%%%%%%%%%%%%%%%%%%%%%%%%%%%%

\section{Organizzazione della tesi}
\label{sec:intro4}

Vi spiego brevemente quali sono le parti di questo template:

\begin{description}
  \item[Susanna.tex] Questo è il file principale del template: contiene le impostazioni generali e la struttura della tesi. Ricordate di impostarlo come documento master ogni volta che iniziate a lavorare alla tesi (ovviamente potete rinominarlo, nabbazzi)
  \item[frontespizio.tex] Sarebbe la copertina della tesi, nonchè la prima pagina. Contiene nome dell'università, del dipartimento, nome della tesi bla bla bla... io ho scelto un argomento di Fisica molto importante <3
  \item[dedica.tex] Una pagina dove scrivete a chi dedicate la tesi (Susanna <3)
  \item[introduzione.tex] Il file che contiene questo capitolo introduttivo (vi consiglio di creare appunto un file .tex per ogni capitolo). Le 4 sezioni di questo capitolo (\nameref{sec:intro1}, \nameref{sec:intro2}, \nameref{sec:intro3} e \nameref{sec:intro4}) sono le 4 sezioni standard da inserire nell'introduzione di una tesi, quindi vi consiglio di lasciarle così
  \item[start.tex] Contiene l'unico capitolo utile di questo documento: spiega come scaricare l'occorrente e come configurare il tutto per lavorare con \LaTeX
  \item[conclusione.tex] Contiene la conclusione (YOU DON'T SAY)
  \item[bibliografia.bib] Contiene i dati relativi alle fonti che citerete nella tesi (ad esempio, ora sto citando un libro sugli algoritmi genetici \cite{5635176}, l'unico inserito nella bibliografia di questo template)
  \item[IEEEtran.bst] È lo stile della bibliografia, non lo toccate
  \item[imgs] Cartella contenente le immagini (YOU DON'T SAY AGAIN)
\end{description}
	%%%%%%%%%%%%%%%%%%%%%%%%%%%%%%%%%%%%%%%%%%%%%%%%%%%%%%%%%%%

\chapter{Come fare le cose}
\label{chap:omfglol}

In questo capitolo vediamo la roba smanettosa per iniziare a smanettare

%%%%%%%%%%%%%%%%%%%%%%%%%%%%%%%%%%%%%%%%%%%%%%%%%%%%%%%%%%%

\section{Occorrente}
\label{sec:occorrente}

Roba da scaricare e installare (Tabella \ref{tab:filesize}).

\begin{table}[h]
\centering
\caption{Tabella vergognosamente inutile}
\vspace{0.3cm}
\label{tab:filesize}
\begin{tabular}{lll}
 \hline
File & Piattaforma & Dimensioni \\\hline
\emph{TexMaker} & Windows & 46.3 MB \\
\emph{TexMaker} & Mac & 40.7 MB \\
\emph{MiKTeX} & Windows & 154.1 MB \\
\emph{MacTex} & Mac & 2.3 GB \\
Template Susanna & Multiglobale-powa & 4 MB (circa) \\\hline
\end{tabular}
\end{table}

\pagebreak

%%%%%%%%%%%%%%%%%%%%%%%%%%%%%%%%%%%%%%%%%%%%%%%%%%%%%%%%%%%

\subsection{L'IDE}
\label{sub:ide}

Allora... per prima cosa vi serve un IDE, ovvero un programma che vi funga da editor e compilatore (in realtà il compilatore si scarica a parte ma vabbè). Ce ne sono molti in giro ma io vi consiglio \emph{TexMaker}\footnote{http://www.xm1math.net/texmaker/} per due motivi:

\begin{enumerate}
\item È molto intuitivo è ben fatto
\item Esiste sia per Mac che per Windows
\end{enumerate}

Non dovreste avere problemi con il download e l'installazione (vi state per laureare porca paletta, non devo spiegarvi pure questo).

%%%%%%%%%%%%%%%%%%%%%%%%%%%%%%%%%%%%%%%%%%%%%%%%%%%%%%%%%%%

\subsection{Il compilatore}
\label{sub:compilatore}

Per quanto riguarda il compilatore il discorso è un po' più complicato. Armatevi di pazienza e scaricate \emph{MiKTeX}\footnote{http://miktex.org/download} se avete Windows oppure \emph{MacTex}\footnote{http://mirror.ctan.org/systems/mac/mactex/MacTeX.pkg} se avete un Mac (mi dispiace ma non conosco un compilatore LaTeX per Linux... se lo trovato fatemelo sapere che aggiorno la guida). Entrambi questi programmi inglobano un ambiente di sviluppo \LaTeX costituito da diversi compilatori che il nostro IDE riconoscerà automaticamente.

\textbf{P.S.} Prima che cominciate a strapparvi i capelli, sì, \emph{MacTex} occupa più di 2 GB... questo perchè comprende tutti i pacchetti necessari per \LaTeX, mentre \emph{MiKTeX} (che occupa solo 150 mb) li scarica volta per volta.

%%%%%%%%%%%%%%%%%%%%%%%%%%%%%%%%%%%%%%%%%%%%%%%%%%%%%%%%%%%

\subsection{Il template}
\label{sub:template}

Trovate il sorgente di questo template ad un link dropbox non meglio specificato\footnote{https://www.dropbox.com/sh/1f06sd7eprongvl/dKsfID1Kwc}

%%%%%%%%%%%%%%%%%%%%%%%%%%%%%%%%%%%%%%%%%%%%%%%%%%%%%%%%%%%

\section{Configurazione dell'IDE}
\label{sec:configurazioneIDE}

Oooh, ora che avete installato IDE e compilatore, lanciate l'IDE. Principalmente dovete fare tre cose una volta avviato:

\begin{enumerate}
\item Aprite il file Susanna.tex del template
\item Andate su Opzioni -> Definire il documento corrente come Master (questo serve per dire all'IDE che gli altri documenti sono inclusi in un documento master e che quindi, al momento della compilazione, non devono essere trattati come documenti separati)
\item Andate nelle preferenze dell'IDE nella sezione Compilazione Rapida e personalizzate la compilazione tramite l'assistente-wizard. Essenzialmente dovete configurarla in modo da effettuare 3 compilazioni: PdfLatex, BibTex e di nuovo PdfLatex. Oltre a queste tre compilazioni aggiungete una quarta opzione ovvero la visualizzazione pdf.
\end{enumerate}

Vi spiego meglio il punto 3... praticamente ci sono più compilatori diversi, che svolgono operazioni diverse... ma a noi interessano solo due compilatori, ovvero PdfLatex (che compila il codice \LaTeX in un documento pdf) e BibTex (che compila la bibliografia). Le compilazioni sono 3 e in quel preciso ordine perchè altrimenti la bibliografia non viene compilata bene (non chiedetemi perchè). Per evitare di dover eseguire manualmente le diverse compilazioni, \emph{TexMaker} vi dà la possibilità di utilizzare la Compilazione Rapida che esegue automaticamente queste operazioni con un click. Configuratela come vi ho spiegato e non avrete problemi.

%%%%%%%%%%%%%%%%%%%%%%%%%%%%%%%%%%%%%%%%%%%%%%%%%%%%%%%%%%%

\section{Siamo pronti}
\label{sec:pronti}

Abbiamo configurato l'IDE ed il (i) compilatore(i). Ora premendo sulla freccina della Compilazione Rapida (oppure premendo F1) potrete compilare il vostro codice \LaTeX in pdf. Fate una prova compilando il template (il pdf purtroppo, così come tutti gli scarti della compilazione, verranno generati nella stessa cartella del sorgente...).

E ora? Ora create i vostri capitoli copiando la struttura di \emph{start.tex} e di \emph{introduzione.tex} ed integrateli nel documento master :) se avete bisogno di ulteriori dettagli sulla sintassi \LaTeX vi consiglio di farvi un giro nella sezione Tex di \emph{Stack Exchange}\footnote{http://tex.stackexchange.com/}: è tipo Yahoo Answers ma focalizzato ovviamente su \LaTeX :)
%	\input{HW1}
%	\input{HW2}
%	\input{HW3}
%	\input{HW4}
%	\input{HW5}
%	\input{HW6}
%	\input{HW7}
	%\input{conclusione}
	% === Bibliografia ====================================
	\newpage
	\bibliographystyle{IEEEtran}
	\bibliography{bibliografia-tesi.bib}
\end{document}
%========================================
